\maketitle

\begin{abstract}
Given the probability measure produced by Paper~I and the dynamics of Paper~II,
this paper asks when an observer can recover the full state space from
time-delayed measurements alone.  The answer is the \emph{reconstruction
theorem}: for a measure-preserving system $(X, \mathcal{B}, \mu, T)$ and a
bounded observation function $h \in L^\infty(X, \mu)$, the following are
equivalent: (i)~the observable $\sigma$-algebra
$\mathcal{O}_h = \sigma(\{h \circ T^n : n \in \mathbb{Z}\})$ equals
$\mathcal{B}(X)$ modulo $\mu$; (ii)~the algebra generated by the delay
observations $\{h \circ T^n\}$ is dense in $L^2(X, \mu)$; (iii)~the delay
map $\Phi_h : x \mapsto (h(T^n x))_{n \in \mathbb{Z}}$ is a
measure-theoretic embedding $X \hookrightarrow \mathbb{R}^\mathbb{Z}$.

The equivalence of~(i) and~(ii) is the \emph{density bridge}: the
$L^2$-closure of an algebra of bounded measurable functions equals the $L^2$
space of the $\sigma$-algebra it generates.  This is proved via the functional
monotone class theorem and simple-function approximation.

The cyclic vector condition from Paper~II — that
$\overline{\mathrm{span}}\{U_T^n h\} = L^2(X,\mu)$ — implies (i)--(iii) but
is strictly stronger: it is a spectral condition on $U_T$, while reconstruction
is a condition on the Boolean algebra of observable events.  The two coincide
when $U_T$ has simple $L^2$-spectrum.

The reconstruction theorem closes the loop with Paper~I: the Stone space
$\mathrm{St}(\mathcal{O}_h)$ constructed there as a technical tool for measure
extension is, under condition~(i), measure-theoretically isomorphic to $X$
itself.  The programme begins with a compact completion of the sample space and
ends by identifying it with the state space.

Compared to the classical Takens embedding theorem~\cite{Takens1981}, the
present result requires only measurability, uses all finite delays rather than
a fixed window, and replaces the dimension count with the algebraic density
condition.
\end{abstract}

% ====================================================================
\section{Introduction}
\label{sec:intro}

\paragraph{The problem.}
Papers~I and~II established that coherent structured observations determine a
probability measure, and that measure determines the dynamics.  The remaining
question is whether the observer can recover the state space itself.  Given
only the time series $h(x), h(Tx), h(T^2 x), \ldots$, can they reconstruct
where in $X$ they are?

The answer depends on how much information the observation function $h$ carries
about the state.  If $h$ is constant, no reconstruction is possible.  If the
orbit $\{h \circ T^n : n \in \mathbb{Z}\}$ collectively separates almost every
pair of points, reconstruction holds.  The reconstruction theorem makes this
precise.

\paragraph{The density bridge.}
The central result is that two conditions — one on the Boolean algebra of
observable events, one on the $L^2$-approximation theory of the delay
observations — are equivalent.  The bridge between them is a general lemma on
function algebras: for a subalgebra
$\mathcal{A} \subseteq L^\infty(X, \mu)$, the $L^2$-closure of $\mathcal{A}$
equals $L^2(X, \sigma(\mathcal{A}), \mu)$.  This is an instance of the
functional monotone class theorem.

\paragraph{The cyclic vector condition.}
The Koopman operator $U_T$ on $L^2(X,\mu)$ provides a spectral perspective.
The condition that $h$ is a \emph{cyclic vector} for $U_T$ —
$\overline{\mathrm{span}}\{U_T^n h : n \in \mathbb{Z}\} = L^2$ — implies
reconstruction.  The implication is strict: reconstruction asks only that the
\emph{algebra} generated by $\{h \circ T^n\}$ is $L^2$-dense, while cyclicity
asks that the \emph{linear span} is dense.  These coincide if and only if
$U_T$ has simple $L^2$-spectrum.

\paragraph{The Stone space connection.}
Paper~I~\cite{mahon_paper1} built the Stone space
$\mathrm{St}(\mathcal{C})$ as a compact Hausdorff completion of the sample
space for the purpose of measure extension.  Paper~III identifies the
Stone space of the observable algebra $\mathcal{O}_h$ with the state space $X$
when reconstruction holds.  The compact space introduced as a technical device
in Paper~I reappears here as the object being recovered.

\paragraph{Relation to Takens.}
The classical Takens theorem~\cite{Takens1981} requires a smooth vector field
on a compact $d$-manifold, a generic observation function, and $2d+1$ delays;
it produces a diffeomorphic embedding.  The reconstruction theorem here requires
only measurability, uses all finite delays, and replaces the dimension count
with the algebraic density condition.  On a compact manifold, the density
condition holds generically, recovering the spirit of Takens without
differentiability hypotheses.

\paragraph{Organisation.}
Section~\ref{sec:setup} sets up the dynamical system, observation function, and
observable algebra.
Section~\ref{sec:density-bridge} proves the density bridge lemma.
Section~\ref{sec:reconstruction} states and proves the reconstruction theorem.
Section~\ref{sec:cyclic} places the cyclic vector condition in context.
Section~\ref{sec:stone} closes the loop with Paper~I.
Section~\ref{sec:takens} compares with the classical Takens theorem.

% ====================================================================
\section{Setup}
\label{sec:setup}

Throughout this paper, $(X, \mathcal{B}, \mu)$ is a standard probability space
and $T : X \to X$ is a measurable, $\mu$-preserving map.  We assume $T$ is
\emph{invertible}: $T$ is a bimeasurable bijection with $T_*\mu = \mu$ and
$(T^{-1})_*\mu = \mu$.

We fix a bounded observation function $h \in L^\infty(X, \mu)$.  The
invertibility of $T$ allows us to use both forward and backward delays:

\begin{definition}[Observable algebra]
\label{def:obs-algebra}
The \emph{observable algebra} generated by $h$ is
\[
  \mathcal{O}_h \;:=\; \sigma\!\bigl(\{h \circ T^n : n \in \mathbb{Z}\}\bigr)
  \;\subseteq\; \mathcal{B}(X).
\]
The \emph{delay algebra} is the algebra of finite polynomials in the delay
observations:
\[
  \mathcal{A}_h
  \;:=\;
  \alg\!\bigl(\{h \circ T^n : n \in \mathbb{Z}\} \cup \{1\}\bigr)
  \;\subseteq\; L^\infty(X, \mu),
\]
the smallest subalgebra of $L^\infty(X, \mu)$ containing $1$ and every
$h \circ T^n$.  Concretely, $\mathcal{A}_h$ consists of all finite sums
$\sum_\alpha c_\alpha \prod_{k} (h \circ T^{n_k})^{e_k}$ with
$c_\alpha \in \mathbb{R}$, $n_k \in \mathbb{Z}$, $e_k \in \mathbb{N}$.
\end{definition}

\begin{definition}[Delay map]
\label{def:delay-map}
The \emph{delay map} is
\[
  \Phi_h : X \to \mathbb{R}^\mathbb{Z},
  \qquad
  \Phi_h(x) \;=\; \bigl(h(T^n x)\bigr)_{n \in \mathbb{Z}}.
\]
Since each coordinate $x \mapsto h(T^n x)$ is measurable and $\mathbb{R}^\mathbb{Z}$
carries the product Borel $\sigma$-algebra $\mathcal{B}(\mathbb{R}^\mathbb{Z}) =
\sigma(\{\pi_n : n \in \mathbb{Z}\})$, the map $\Phi_h$ is measurable.
\end{definition}

\begin{lemma}[Observable algebra as pullback]
\label{lem:obs-pullback}
$\mathcal{O}_h = \Phi_h^{-1}(\mathcal{B}(\mathbb{R}^\mathbb{Z}))$ mod $\mu$.
\end{lemma}

\begin{proof}
The product $\sigma$-algebra $\mathcal{B}(\mathbb{R}^\mathbb{Z})$ is generated by
coordinate projections $\{\pi_n : n \in \mathbb{Z}\}$.  Since
$\pi_n \circ \Phi_h = h \circ T^n$, the pullback
$\Phi_h^{-1}(\mathcal{B}(\mathbb{R}^\mathbb{Z}))$ is generated by
$\{(h \circ T^n)^{-1}(E) : E \in \mathcal{B}(\mathbb{R}),\, n \in \mathbb{Z}\}$,
which is exactly $\mathcal{O}_h$.
\end{proof}

\noindent The observation that $\mathcal{O}_h$ is the pullback $\sigma$-algebra of
$\Phi_h$ means that $\mathcal{O}_h = \mathcal{B}$ mod $\mu$ is precisely the condition
that $\Phi_h$ is a measure-theoretic embedding — condition~(iii) of the reconstruction
theorem.  The content of the theorem is that this geometric condition is equivalent to
the analytic condition~(ii) on algebra density.

% ====================================================================
\section{The density bridge}
\label{sec:density-bridge}

The density bridge is the lemma connecting the Boolean algebra level (the
$\sigma$-algebra $\mathcal{O}_h$) to the $L^2$ level (approximation by elements
of $\mathcal{A}_h$).  It holds in complete generality for any subalgebra of
$L^\infty$.

\begin{lemma}[Density bridge]
\label{lem:density-bridge}
Let $(X, \mathcal{B}, \mu)$ be a probability space and let
$\mathcal{A} \subseteq L^\infty(X, \mu)$ be a subalgebra containing the
constant function~$1$.  Then:
\[
  \overline{\mathcal{A}}^{L^2}
  \;=\;
  L^2(X,\, \sigma(\mathcal{A}),\, \mu).
\]
Consequently, $\overline{\mathcal{A}}^{L^2} = L^2(X, \mu)$ if and only if
$\sigma(\mathcal{A}) = \mathcal{B}$ modulo $\mu$.
\end{lemma}

\begin{proof}
Write $\mathcal{M} := \sigma(\mathcal{A})$ and $V := \overline{\mathcal{A}}^{L^2}$.

\textbf{Step 1: $V \subseteq L^2(X, \mathcal{M}, \mu)$.}
Every $f \in \mathcal{A}$ is $\mathcal{A}$-measurable, hence $\mathcal{M}$-measurable.
So $\mathcal{A} \subseteq L^2(X, \mathcal{M}, \mu)$.  Since $L^2(X, \mathcal{M}, \mu)$
is a closed subspace of $L^2(X, \mathcal{B}, \mu)$, taking the closure gives
$V \subseteq L^2(X, \mathcal{M}, \mu)$.

\textbf{Step 2: $L^2(X, \mathcal{M}, \mu) \subseteq V$.}
It suffices to show that every indicator $\mathbf{1}_E$ with $E \in \mathcal{M}$
lies in $V$.  By the functional monotone class theorem~\cite{Dynkin}, the
class
\[
  \mathcal{H}
  \;:=\;
  \bigl\{f \in L^\infty(X, \mathcal{M}, \mu) : f \in V\bigr\}
\]
is a vector space containing $\mathcal{A}$ and closed under bounded monotone
pointwise limits (since if $f_n \in V$ are uniformly bounded and $f_n \to f$
pointwise a.e.\ with $f \in L^\infty$, then $f_n \to f$ in $L^2$ by dominated
convergence, so $f \in V$).  Since $\mathcal{A}$ is an algebra containing $1$
and $\sigma(\mathcal{A}) = \mathcal{M}$, the monotone class theorem gives
$\mathcal{H} = L^\infty(X, \mathcal{M}, \mu)$.  In particular every
$\mathcal{M}$-indicator $\mathbf{1}_E \in V$.

Since the $\mathcal{M}$-simple functions are dense in $L^2(X, \mathcal{M}, \mu)$
and $V$ is closed, $L^2(X, \mathcal{M}, \mu) \subseteq V$.

\textbf{Conclusion.}
Steps~1 and~2 give $V = L^2(X, \mathcal{M}, \mu)$.  The final equivalence is
immediate: $V = L^2(X, \mathcal{B}, \mu)$ iff $L^2(X, \mathcal{M}, \mu) =
L^2(X, \mathcal{B}, \mu)$ iff $\mathcal{M} = \mathcal{B}$ mod $\mu$.
\end{proof}

\begin{remark}[Role of the algebra structure]
\label{rem:algebra-necessary}
The algebra structure of $\mathcal{A}$ is essential.  The corresponding
statement for the closed \emph{linear} span $\mathcal{H}_{\mathcal{A}} =
\overline{\mathrm{span}}(\mathcal{A})$ is false in general:
$\mathcal{H}_{\mathcal{A}} = L^2$ does not imply
$\sigma(\mathcal{A}) = \mathcal{B}$.

For example, take $X = [0,1]$, $\mu = $ Lebesgue, $T(x) = 2x \bmod 1$ (the
doubling map), and $h(x) = x$.  The orbit $\{h \circ T^n\}(x) = \{2^n x \bmod 1\}$
generates $\mathcal{B}([0,1])$ as a $\sigma$-algebra (the dyadic structure separates
points).  But $\mathrm{span}\{2^n x \bmod 1 : n \geq 0\}$ is a proper subspace of
$L^2[0,1]$: it does not contain $x^2$.  So the linear span can be a proper dense
subspace even when $\sigma(\mathcal{A}) = \mathcal{B}$.

In the other direction, $h$ cyclic for $U_T$ (linear span dense) always implies
$\sigma(\mathcal{A}_h) = \mathcal{B}$ mod $\mu$ — see Corollary~\ref{cor:cyclic-implies-reconstruction}.
\end{remark}

% ====================================================================
\section{The reconstruction theorem}
\label{sec:reconstruction}

\begin{theorem}[Reconstruction theorem]
\label{thm:reconstruction}
Let $(X, \mathcal{B}, \mu, T)$ be an invertible measure-preserving system and
$h \in L^\infty(X, \mu)$.  The following are equivalent:
\begin{enumerate}[label=(\roman*)]
  \item $\mathcal{O}_h = \mathcal{B}$ modulo $\mu$;
  \item $\mathcal{A}_h$ is dense in $L^2(X, \mu)$;
  \item $\Phi_h : X \to \mathbb{R}^\mathbb{Z}$ is a measure-theoretic embedding:
    $\Phi_h^{-1}(\mathcal{B}(\mathbb{R}^\mathbb{Z})) = \mathcal{B}$ modulo $\mu$.
\end{enumerate}
\end{theorem}

\begin{proof}
\emph{(i) $\Leftrightarrow$ (ii).}
Apply the density bridge (Lemma~\ref{lem:density-bridge}) with
$\mathcal{A} = \mathcal{A}_h$ and $\sigma(\mathcal{A}_h) = \mathcal{O}_h$:
$\overline{\mathcal{A}_h}^{L^2} = L^2(X, \mu)$ iff $\mathcal{O}_h = \mathcal{B}$
mod $\mu$.

\emph{(i) $\Leftrightarrow$ (iii).}
By Lemma~\ref{lem:obs-pullback},
$\mathcal{O}_h = \Phi_h^{-1}(\mathcal{B}(\mathbb{R}^\mathbb{Z}))$ mod $\mu$.
So $\mathcal{O}_h = \mathcal{B}$ iff
$\Phi_h^{-1}(\mathcal{B}(\mathbb{R}^\mathbb{Z})) = \mathcal{B}$ mod $\mu$,
which is exactly condition~(iii).
\end{proof}

\begin{remark}[What ``measure-theoretic embedding'' means]
\label{rem:embedding}
Condition~(iii) says that $\Phi_h$ is \emph{almost everywhere injective} and
the pullback $\sigma$-algebra is all of $\mathcal{B}$.  More precisely: since
$\Phi_h^{-1}(\mathcal{B}(\mathbb{R}^\mathbb{Z})) = \mathcal{B}$ mod $\mu$, for
any two sets $A, A' \in \mathcal{B}$ with $\mu(A \triangle A') > 0$ there exists
a Borel set $E \subseteq \mathbb{R}^\mathbb{Z}$ with
$\mu(\Phi_h^{-1}(E) \triangle A) = 0$ and
$\mu(\Phi_h^{-1}(E) \triangle A') > 0$.  In measure-theoretic language,
$\Phi_h$ induces an isomorphism of measure algebras
$(X/\mu, \mathcal{B}/\mu) \cong (\mathbb{R}^\mathbb{Z}/(\Phi_h)_*\mu,
\mathcal{B}(\mathbb{R}^\mathbb{Z})/(\Phi_h)_*\mu)$.

The pushforward measure $(\Phi_h)_*\mu$ is concentrated on the image
$\Phi_h(X) \subseteq \mathbb{R}^\mathbb{Z}$ by definition.  The
$T$-invariance of $\mu$ implies that $(\Phi_h)_*\mu$ is invariant under the
bilateral shift $\sigma$ on $\mathbb{R}^\mathbb{Z}$: the dynamical system
$(X, T, \mu)$ is isomorphic to $(\mathbb{R}^\mathbb{Z}, \sigma, (\Phi_h)_*\mu)$
restricted to the image of $\Phi_h$.
\end{remark}

\begin{corollary}[Shift invariance]
\label{cor:shift-invariance}
Under condition~(i)--(iii), the map $\Phi_h$ intertwines $T$ on $X$ with the
bilateral shift $\sigma$ on $\mathbb{R}^\mathbb{Z}$: $\Phi_h \circ T = \sigma
\circ \Phi_h$ $\mu$-almost everywhere.
\end{corollary}

\begin{proof}
For each $n \in \mathbb{Z}$,
$(\Phi_h(Tx))_n = h(T^n(Tx)) = h(T^{n+1}x) = (\Phi_h(x))_{n+1} =
(\sigma(\Phi_h(x)))_n$.
\end{proof}

% ====================================================================
\section{The cyclic vector condition}
\label{sec:cyclic}

The Koopman operator $U_T$ on $L^2(X, \mu)$ encodes the dynamics at the
operator level.  Recall from Paper~II~\cite{mahon_paper2} that $U_T$ is unitary
(since $T$ is invertible and measure-preserving) and that the predictive
operators $\{K_t\}$ are related to $U_T$ via the Koopman--Perron duality.

\begin{definition}[Cyclic vector]
\label{def:cyclic}
The function $h \in L^2(X, \mu)$ is a \emph{cyclic vector} for $U_T$ if
\[
  \mathcal{H}_h
  \;:=\;
  \overline{\mathrm{span}}\!\bigl\{U_T^n h : n \in \mathbb{Z}\bigr\}
  \;=\; L^2(X, \mu).
\]
\end{definition}

\begin{corollary}[Cyclic implies reconstruction]
\label{cor:cyclic-implies-reconstruction}
If $h \in L^\infty(X, \mu)$ is a cyclic vector for $U_T$, then
$\mathcal{O}_h = \mathcal{B}$ modulo $\mu$.
\end{corollary}

\begin{proof}
Since $\mathcal{A}_h$ contains $\mathrm{span}\{h \circ T^n\} =
\mathrm{span}\{U_T^n h\}$ as a subset, $\overline{\mathcal{A}_h}^{L^2}
\supseteq \mathcal{H}_h = L^2$.  By the density bridge, $\mathcal{O}_h =
\mathcal{B}$ mod $\mu$.
\end{proof}

\begin{remark}[The implication is strict]
\label{rem:strict}
The converse of Corollary~\ref{cor:cyclic-implies-reconstruction} fails in
general: $\mathcal{O}_h = \mathcal{B}$ does not imply $h$ is cyclic.
The doubling map example of Remark~\ref{rem:algebra-necessary} witnesses this.

The cyclic vector condition asks for density of the \emph{linear span} of the
orbit; reconstruction asks only for density of the \emph{algebra}.  For a bounded
$h$, the linear span is contained in the algebra, so cyclicity is the stronger
condition.

The two conditions coincide precisely when $U_T$ has \emph{simple
$L^2$-spectrum}: $L^2(X,\mu)$ is a cyclic $U_T$-module, meaning there exists
some cyclic vector.  In this case, every generating observation $h$ (satisfying
$\mathcal{O}_h = \mathcal{B}$) is automatically cyclic.  Simple spectrum is
a generic property of ergodic transformations in the Baire category sense
\cite{Halmos1944} but is not universal.
\end{remark}

\begin{remark}[Operator-algebraic perspective]
\label{rem:vna}
The distinction between reconstruction and cyclicity reflects a level in the
Stone--Gelfand--Naimark duality:

\begin{center}
\begin{tabular}{lll}
\textbf{Level} & \textbf{Object} & \textbf{Reconstruction condition} \\
\hline
Boolean algebra & $\mathcal{O}_h \subseteq \mathcal{B}$ &
  $\mathcal{O}_h = \mathcal{B}$ mod $\mu$ \\
$L^2$ function space & $L^2(\mathcal{O}_h) \subseteq L^2(\mathcal{B})$ &
  $L^2(\mathcal{O}_h) = L^2(\mathcal{B})$ \\
von Neumann algebra & $\{U_T^n h\}'' \subseteq \mathcal{B}(L^2)$ &
  $\{U_T^n h\}'' = L^\infty(X,\mu)$ \\
\end{tabular}
\end{center}

The reconstruction theorem operates at the first two levels, which are
equivalent by the density bridge.  The cyclic vector condition operates at the
third level.  These coincide when $L^\infty(X,\mu)$ is generated as a von
Neumann algebra by the orbit of $h$ under $U_T$ \emph{and} $h$ is a cyclic
vector — which holds iff $U_T$ has simple spectrum and $h$ is a generator.
\end{remark}

% ====================================================================
\section{The Stone space identification}
\label{sec:stone}

Paper~I constructed the Stone space $\mathrm{St}(\mathcal{C})$ of the cylinder
algebra $\mathcal{C}$ as the canonical compact Hausdorff completion of the
sample space $\Omega$.  In the present setting, the cylinder algebra is replaced
by the observable algebra $\mathcal{O}_h$, and the Stone space
$\mathrm{St}(\mathcal{O}_h)$ is the compact completion of $X$ for the topology
generated by observable events.

\begin{proposition}[Stone space identification]
\label{prop:stone-id}
Suppose condition~(i) of Theorem~\ref{thm:reconstruction} holds:
$\mathcal{O}_h = \mathcal{B}$ modulo $\mu$.  Then the Stone embedding
\[
  \iota : X \to \mathrm{St}(\mathcal{O}_h),
  \qquad
  \iota(x) = \{E \in \mathcal{O}_h : x \in E\},
\]
is a measure-theoretic isomorphism: the pushforward $\iota_*\mu$ is a
probability measure on $\mathrm{St}(\mathcal{O}_h)$, and $\iota$ identifies the
measure algebra $(X/\mu, \mathcal{B}/\mu)$ with
$(\mathrm{St}(\mathcal{O}_h)/\iota_*\mu,
\mathcal{B}(\mathrm{St}(\mathcal{O}_h))/\iota_*\mu)$.
\end{proposition}

\begin{proof}
When $\mathcal{O}_h = \mathcal{B}$ mod $\mu$, the Stone embedding $\iota$ has
the following properties.

\emph{Injectivity mod $\mu$:} if $\iota(x) = \iota(x')$ then $x$ and $x'$
belong to the same sets in $\mathcal{O}_h$.  Since $\mathcal{O}_h = \mathcal{B}$
mod $\mu$, any set separating $x$ from $x'$ can be approximated mod $\mu$ by an
$\mathcal{O}_h$-set.  Injectivity follows for $\mu$-almost every pair.

\emph{$\sigma$-algebra identification:} the pullback $\iota^{-1}$ maps the Baire
$\sigma$-algebra of $\mathrm{St}(\mathcal{O}_h)$ (generated by clopen sets
$[E] = \{u : E \in u\}$ for $E \in \mathcal{O}_h$) to $\mathcal{O}_h =
\mathcal{B}$ mod $\mu$.

\emph{Measure recovery:} CE (Collective Exhaustion), which holds since $\mu$ is
$\sigma$-additive on $\mathcal{O}_h = \mathcal{B}$, forces $\iota_*\mu$ to
concentrate on the principal ultrafilters $\iota(X) \subset
\mathrm{St}(\mathcal{O}_h)$ (Paper~I, Proposition~5.5).  So $\iota_*\mu$ is
supported on the image $\iota(X)$ and the measure algebras are identified.
\end{proof}

\begin{remark}[The programme closes]
The Stone space $\mathrm{St}(\mathcal{C})$ introduced in Paper~I as a tool for
deriving $\sigma$-additivity from compactness reappears here as the
reconstruction of $X$.  The programme begins with an observer whose distinctions
generate a Boolean algebra $\mathcal{C}$, and ends by showing that when those
distinctions are complete ($\mathcal{O}_h = \mathcal{B}$), the Stone space of
$\mathcal{C}$ is the state space of the system they are observing.

The three papers trace the same duality from different angles: Paper~I from the
Boolean side (Stone duality on the cylinder algebra), Paper~II from the spectral
side (Koopman operator on $L^2$), Paper~III from the algebraic side (density
bridge between the Boolean and $L^2$ levels).
\end{remark}

% ====================================================================
\section{Comparison with Takens}
\label{sec:takens}

The classical Takens embedding theorem~\cite{Takens1981} states: for a smooth
vector field on a compact $d$-dimensional manifold $M$, and for a generic smooth
observation function $h : M \to \mathbb{R}$, the delay map
\[
  \Phi_h^{(k)} : M \to \mathbb{R}^k,
  \qquad
  \Phi_h^{(k)}(x) = (h(x), h(Tx), \ldots, h(T^{k-1}x))
\]
is an embedding (injective immersion) for $k \geq 2d + 1$.

The reconstruction theorem differs in four respects:

\begin{enumerate}[label=(\arabic*)]
\item \emph{Regularity.}  Takens requires a $C^2$ diffeomorphism and a $C^2$
  observation.  The reconstruction theorem requires only measurability of $T$
  and $h$.

\item \emph{Window length.}  Takens uses a finite window of length $2d+1$,
  determined by the dimension of $M$.  The reconstruction theorem uses all
  delays $n \in \mathbb{Z}$ simultaneously.  The dimension bound in Takens
  ensures that the finite-delay algebra already separates points; here the
  algebra density condition replaces it.

\item \emph{Embedding type.}  Takens produces a diffeomorphic embedding into
  $\mathbb{R}^{2d+1}$.  The reconstruction theorem produces a measure-theoretic
  embedding into $\mathbb{R}^\mathbb{Z}$: an isomorphism of measure algebras,
  not a topological embedding.

\item \emph{Genericity.}  Takens requires genericity of $h$ (to avoid
  degenerate observation functions).  In the measure-theoretic setting, the
  corresponding condition is $\mathcal{O}_h = \mathcal{B}$: the observation
  function generates the full $\sigma$-algebra.  On a compact manifold, this
  holds for any $h$ whose level sets $\{h = c\}$ are $(d-1)$-dimensional
  submanifolds, which is a generic condition on $h$.
\end{enumerate}

\begin{remark}[Recovering Takens on manifolds]
On a compact $d$-manifold with a $C^2$ diffeomorphism $T$ and a generic smooth
$h$, the orbit $\{h \circ T^n : n \in \mathbb{Z}\}$ separates points.  By a
Stone--Weierstrass argument, the algebra $\mathcal{A}_h$ is uniformly dense in
$C(M)$ and hence $L^2$-dense.  So the algebraic density condition of
Theorem~\ref{thm:reconstruction}(ii) holds, and $\Phi_h$ is a
measure-theoretic embedding.  For Lebesgue measure on $M$, this is the
measure-theoretic content of the Takens theorem.
\end{remark}

% ====================================================================
\bibliographystyle{plain}
\bibliography{references}
